\documentclass[12pt, a4paper]{article}

% --- Pacotes Básicos ---
\usepackage[utf8]{inputenc}
\usepackage[]{graphicx}
\usepackage[T1]{fontenc}
\usepackage{times}
\usepackage{setspace}
\onehalfspacing % Set line spacing to 1.5
\setlength{\parindent}{1.25cm}
\setlength{\lineskip}{1.5cm}
%\usepackage[main=brazilian, english]{babel}
\usepackage{indentfirst}
\usepackage{geometry}
\geometry{left=3cm, top=3cm, right=2cm, bottom=2cm}

\usepackage{etoolbox} % Pacote necessário para ajustar o espaçamento

% Criando o ambiente 'citacao'
\newenvironment{citacao}{
  \vspace{0.5cm} % Espaço antes da citação
  \small % Letra menor (geralmente tamanho 10)
  \begin{list}{}{
    \leftmargin=4cm % Recuo de 4 cm da margem esquerda
    \rightmargin=0cm
    \parsep=0pt
    \topsep=0pt
    \partopsep=0pt
  }
  \item\relax
  \linespread{1.0}\selectfont % Garante espaçamento simples dentro da citação
}{
  \end{list}
  \vspace{0.5cm} % Espaço depois da citação
}
\usepackage{fancyhdr}

% Configura o estilo "fancy"
\fancypagestyle{plain}{
  \fancyhf{} % Limpa todos os cabeçalhos e rodapés
  \fancyfoot[R]{\thepage} % Define a paginação no rodapé à direita (R = Right)
  \renewcommand{\headrulewidth}{0pt}  % Remove a linha divisória do cabeçalho
  \renewcommand{\footrulewidth}{0pt}  % Remove a linha divisória do rodapé
}

% Ativa o estilo configurado para todo o documento
\pagestyle{plain}
% --- Informações do Documento ---
\title{UNIVERSIDADE FEDERAL DA FRONTEIRA SUL\\Ciência da Computação\\Meio Ambiente, Economia e Sociedade\\Prof. Leandro Bordin\\ Movimento Ambientalista: Desenvolvimento Sustentável em Meio ao Capitalismo}
\date{}

\begin{document}
\maketitle

\begin{flushright}
	Erickson Giesel Muller
\end{flushright}

Três correntes da economia ambiental: economia ambiental neoclássica, economia ambiental ecológica e economia ambiental marxista.












%\begin{thebibliography}{99}
%
%\bibitem{curta2012}
%
%\bibitem{mabilde1983} 
%	MABILDE, Pierre F. A. B. \textbf{Apontamentos sobre os indígenas selvagens da Nação Coroados das matas da Província do Rio Grande do Sul}: 1836-1866. São Paulo: IBRASA, Fundação Nacional Pró-Memória, 1983.
%
%\bibitem{prous1992}
%	PROUS, André. \textbf{Arqueologia Brasileira}. Brasília: Editora da UnB, 1992.
%
%\bibitem{tommasino1995}
%	TOMMASINO,  Kimiye.  \textbf{A  História  dos  Kaingáng  da  Bacia  do  Tibagi:  Uma  Sociedade Jê Meridional em Movimento}. São Paulo: USP, 1995. Tese (Doutorado  em Antropologia) – Instituto de Filosofia e Ciências Humanas, Universidade de São  Paulo, 1995. 
%\end{thebibliography}
%
\end{document}
